\chapter{The Conditional Critical-Axis Theorem}
\label{chap:conditional}

\section{Statement of the theorem}

The central conditional result is deliberately short. Its purpose is to expose exactly where the Riemann Hypothesis enters.

\begin{theorem}[SOH-T001: conditional critical-axis theorem]
\label{thm:conditional-rh}
Suppose there exists a state map $s\mapsto\ket{\Psi(s)}$ on the open critical strip satisfying:
\begin{enumerate}
  \item for $s=\sigma+it$, the normalized channel weights are $\sigma$ and $1-\sigma$;
  \item the zero readout is the equal-gain amplitude $R(s)=\braket{+}{\Psi(s)}$;
  \item $\xi(s)=0$ if and only if $R(s)=0$.
\end{enumerate}
Then every non-trivial zero of $\zeta(s)$ has real part $1/2$.
\end{theorem}
\begin{proof}
Let $\rho=\beta+i\gamma$ be a non-trivial zero. By the third hypothesis, $R(\rho)=0$. By the first two hypotheses, $R$ is proportional to
\[
 \sqrt\beta+e^{i\Phi(\rho)}\sqrt{1-\beta}.
\]
The exact cancellation lemma \cref{lem:exact-cancellation} implies $\beta=1/2$. Since $\rho$ was arbitrary, all non-trivial zeros lie on the critical line.
\end{proof}

\status{Conditional theorem. The proof is exact, but hypothesis 3 is SOH-C001 and is not established. Equal-gain normalization also requires canonical justification.}

\section{Equivalent positive form}

The theorem can be stated without choosing a phase representative.

\begin{corollary}[Positive detector formulation]
Suppose there are a normalized state $\Psi(s)$ and a positive function $w(s)$ on $\Sstrip$ such that
\begin{equation}
 |\xi(s)|^2=w(s)\bra{\Psi(s)}P_+\ket{\Psi(s)},
 \label{eq:conditional-positive}
\end{equation}
where the component norms of $\Psi(s)$ are $\sqrt\sigma$ and $\sqrt{1-\sigma}$. Then RH holds.
\end{corollary}
\begin{proof}
At a zero of $\xi$, positivity of $w$ and $P_+$ implies $P_+\Psi=0$. The cancellation lemma gives $\sigma=1/2$.
\end{proof}

This formulation is a candidate target for a state-map kernel construction because $|\xi|^2$ and Hilbert-space norms are non-holomorphic but positive. No identity of the displayed form is assumed elsewhere unless explicitly stated.

\section{Eta version}

By \cref{prop:eta-zeta-equivalence}, the completed zeta factor can be replaced by eta in the open strip.

\begin{corollary}[Conditional eta theorem]
If a canonical normalized binary state has equal-gain readout zero exactly when $\eta(s)=0$ for $s\in\Sstrip$, then RH holds.
\end{corollary}

The eta formulation uses the exact alternating-series zero equivalence; the xi formulation uses the functional symmetry and entire completion. A successful state-map construction would have to specify which analytic object supplies each required property.

\section{Why the theorem is not tautological}

The conclusion follows immediately once zerohood is represented by normalized two-channel cancellation. This may appear to make the theorem trivial. The same logical form occurs in spectral reformulations: if the relevant zeros are eigenvalues of a self-adjoint operator, their spectral parameters are real. In both cases the non-trivial content lies in constructing the object with all required equivalences.

The theorem becomes tautological if the state is defined from the zero condition. The following construction is invalid as an explanation:
\[
 \ket{\Psi(s)}=
 \begin{cases}
  \ket{-},&\xi(s)=0,\\
  \ket{+},&\xi(s)\neq0.
 \end{cases}
\]
It satisfies zero equivalence but is discontinuous, non-canonical, and uses the answer as input. Even a smooth interpolation built from $|\xi(s)|$ may merely repackage the zero set.

The meaningful problem is therefore not existence in the set-theoretic sense. It is existence with the stated arithmetic, analytic, symmetry, metric, and positivity structure.

\section{Multiplicity and local form}

Assume the stronger factorization
\[
 \xi(s)=g(s)R(s)
\]
with $g$ non-zero and sufficiently regular. Then the order of a zero of $\xi$ equals the order of the readout zero. Near a simple zero $\rho$,
\[
 R(s)=c(s-\rho)+O((s-\rho)^2),\qquad c\neq0.
\]
But the normalized amplitude $R(s)$ depends on both $s$ and $\bar s$ through $\sigma=\Re s$. It is therefore not holomorphic in the naive parameterization. A rigorous local multiplicity statement must specify whether $R$ is a complex-analytic scalar, a real-analytic section, or a norm of an analytic vector. This technical issue is not cosmetic: holomorphic zero multiplicity is central to zeta theory.

One construction target is to let an analytic vector $v(s)$ carry the holomorphic structure and then derive the normalized state
\[
 \Psi(s)=\frac{v(s)}{\|v(s)\|}
\]
only away from its singular set. The detector amplitude of $v(s)$ can be holomorphic even though the normalized probabilities are not. A non-zero normalization factor would then preserve zeros. Such candidates are considered in \cref{chap:candidates}.

\section{Necessity of equal channel metric}

The conditional theorem is unstable under detector asymmetry. If the readout is
\[
 R_{a,b}(s)=a\sqrt\sigma+b e^{i\Phi(s)}\sqrt{1-\sigma},
\]
then a zero forces
\[
 \sigma=\frac{b^2}{a^2+b^2},
\]
not necessarily $1/2$. Therefore a proof must establish $a=b$ from a specified symmetry or geometry. For the state-map programme, equal treatment of complementary channels must be encoded in the Hilbert-space metric and detector rather than inferred from the zeta functional symmetry alone.

A corresponding structural requirement is
\begin{equation}
 X^\dagger P_+X=P_+,
 \label{eq:detector-swap-invariance}
\end{equation}
and that the channel basis is orthonormal in a swap-invariant metric. This condition makes equal gain an explicit hypothesis to be derived rather than chosen silently.

\section{Minimal bridge versus explanatory bridge}

The minimal bridge needed for the theorem is zero equivalence. A stronger explanatory construction would additionally:
\begin{itemize}
  \item derive the state from the theta or eta kernel;
  \item determine the zero ordinates without fitting them;
  \item retain the Euler product or explicit-formula arithmetic;
  \item identify an independently proved positive or self-adjoint mechanism excluding horizontal displacement;
  \item reproduce known zero statistics as output rather than input.
\end{itemize}
These stronger goals are not required for the logical implication but are required for a non-circular construction.

\section{The open theorem-construction pair}

For this initial state-map branch the status can be expressed in two lines:
\begin{align*}
 \textbf{Theorem:}&\quad \text{canonical normalized cancellation}\Longrightarrow\text{RH},\\
 \textbf{Open construction:}&\quad \xi\text{ or }\eta\Longrightarrow\text{canonical normalized cancellation}.
\end{align*}
Later quotient and kernel chapters pursue different exact objects and do not automatically instantiate the open second line.
